\chapter{Hepatitis B}
\index{Hepatitis B}

\section*{Definition and Overview}
Hepatitis B is a viral infection of the liver caused by the hepatitis B virus, spread through contact with blood and body fluids. It causes acute hepatitis, and in some people, especially those infected in infancy, the infection becomes chronic, persisting for life and carrying risks of cirrhosis and liver cancer. Safe and highly effective vaccines prevent hepatitis B, and antiviral treatment controls chronic infection and greatly reduces complications. some countries have committed to eliminating hepatitis B as a public health threat.

\section*{Epidemiology}
An estimated 200,000 people live with chronic hepatitis B, most born in countries where the virus is common. Worldwide, around 250--300 million people are chronically infected, with most deaths from cirrhosis and liver cancer. In many countries, the birth dose of hepatitis B vaccine and routine childhood vaccination have dramatically reduced new infections. Universal vaccination, screening of at-risk groups, and effective antiviral treatment mean that elimination of hepatitis B is achievable.

\section*{Aetiology and Pathophysiology}
\begin{itemize}
    \item \textbf{Transmission:} blood and body fluids, including mother-to-child during birth, unprotected sex, sharing needles, and contaminated equipment
    \item \textbf{Perinatal transmission:} the main route in high-prevalence countries; infection in infancy almost always becomes chronic
    \item \textbf{Blood-borne spread:} injecting drug use, tattoos and piercings with unsterile equipment, and needle-stick injuries
    \item \textbf{Pathophysiology:} the virus infects liver cells; the immune response causes liver inflammation, and persistent infection drives progressive fibrosis
    \item \textbf{Chronicity:} the immune system clears infection in most adults but not in babies and young children
    \item \textbf{Complications:} chronic inflammation leads to cirrhosis and hepatocellular carcinoma
    \item \textbf{Immunity:} vaccination or resolved infection produces protective antibodies
\end{itemize}

\section*{Clinical Features}
\begin{itemize}
    \item Most acute infections in adults are silent or cause mild, flu-like symptoms
    \item Acute hepatitis: fatigue, nausea, poor appetite, abdominal pain, dark urine, pale stools, and jaundice
    \item Fulminant hepatitis is rare but life-threatening
    \item Chronic hepatitis B is usually asymptomatic for decades
    \item Fatigue may be the only symptom of chronic infection
    \item Advanced disease: weight loss, abdominal swelling, bruising, bleeding, and confusion
    \item Liver cancer may be the first presentation
\end{itemize}

\section*{Differential Diagnosis}
\begin{itemize}
    \item Other causes of acute hepatitis: hepatitis A, C, and E, medications, and alcohol
    \item Other causes of chronic liver disease: fatty liver, alcohol, autoimmune hepatitis, and haemochromatosis
    \item The specific hepatitis B serology makes the diagnosis
    \item Fatigue from other causes is common and should not be attributed to hepatitis B without testing
    \item Jaundice has many causes, including gallstones and pancreatic disease
\end{itemize}

\section*{When to Seek Medical Care}
Testing for hepatitis B is recommended for pregnant women, people born in countries with high prevalence, household and sexual contacts of infected people, people who inject drugs, men who have sex with men, and people with abnormal liver tests. Anyone with symptoms of hepatitis should be tested.

\textbf{Contact your local emergency services for:}
\begin{itemize}
    \item Confusion, drowsiness, or personality change with liver disease
    \item Vomiting blood or passing black stools
    \item Severe abdominal pain or swelling
    \item Sudden severe jaundice with bleeding
\end{itemize}

\section*{Diagnosis and Investigations}
\begin{itemize}
    \item \textbf{Hepatitis B surface antigen:} confirms current infection
    \item \textbf{Surface and core antibodies:} distinguish acute, chronic, and resolved infection and vaccination
    \item \textbf{Viral DNA (viral load):} measures viral activity and guides treatment
    \item \textbf{Liver function tests:} monitor liver inflammation
    \item \textbf{Liver stiffness measurement and ultrasound:} assess fibrosis and screen for liver cancer
    \item \textbf{Liver cancer surveillance:} six-monthly scans and alpha-fetoprotein testing in people at risk
    \item \textbf{Testing of contacts:} household and sexual partners are tested and offered vaccination
\end{itemize}

\section*{Management}
\begin{itemize}
    \item \textbf{Acute hepatitis B:} supportive care; most adults recover completely
    \item \textbf{Chronic hepatitis B:} regular monitoring, with antiviral tablets (such as tenofovir or entecavir) when the virus is active or liver damage is developing
    \item \textbf{Antiviral therapy:} suppresses the virus, heals liver inflammation, and prevents cirrhosis and cancer; usually taken long term
    \item \textbf{Babies of infected mothers:} vaccinated and given hepatitis B immunoglobulin at birth, which prevents infection in more than 95\% of cases
    \item \textbf{Surveillance:} lifelong monitoring for liver cancer in chronic infection
    \item \textbf{Lifestyle:} limit alcohol and maintain a healthy weight
    \item \textbf{Liver cancer and transplant care:} specialist treatment for advanced disease
\end{itemize}

\section*{Complications}
\begin{itemize}
    \item Chronic infection in up to 90\% of infants and 5--10\% of adults
    \item Cirrhosis and liver failure
    \item Hepatocellular carcinoma, the leading cause of liver cancer deaths
    \item Fulminant hepatitis in rare acute cases
    \item Transmission to partners and children
    \item Stigma and psychological impact of diagnosis
\end{itemize}

\section*{Prognosis and Outcomes}
The outlook for hepatitis B has improved enormously. Vaccination prevents infection, and antiviral therapy controls chronic disease, reducing the risk of cirrhosis and liver cancer by more than half. Most people with chronic hepatitis B on treatment live normal lives with near-normal life expectancy. With six-monthly surveillance, liver cancers are detected early and treated effectively. Elimination of hepatitis B is a realistic goal with universal vaccination, screening, and treatment.

\section*{Prevention and Self-Care}
\begin{itemize}
    \item Vaccination: hepatitis B vaccine is part of the childhood schedule and is recommended for all adults at risk
    \item Babies receive the vaccine at birth
    \item Use condoms and avoid sharing needles, syringes, razors, and toothbrushes
    \item Ensure tattoos and piercings are done with sterile equipment
    \item Health workers use gloves and safe needle practices
    \item People with chronic hepatitis B attend regular specialist review
    \item Pregnant women are tested and babies protected at birth
    \item Limit alcohol and keep a healthy weight
\end{itemize}

\section*{Sources and Further Reading}
The following reputable sources provide further information for healthcare students and patients seeking reliable, current advice about hepatitis B.

\begin{itemize}
    \item Hepatitis Australia. \textit{Hepatitis B information and support.} https://www.hepatitisaustralia.com/
    \item Healthdirect. \textit{Hepatitis B: causes, symptoms and treatment.} https://www.healthdirect.gov.au/hepatitis-b
    \item Australian Government Department of Health. \textit{National hepatitis B strategy.} https://www.health.gov.au/
    \item World Health Organization. \textit{Hepatitis B: fact sheet.} https://www.who.int/
    \item NHS. \textit{Hepatitis B: overview and treatment.} https://www.nhs.uk/conditions/hepatitis-b/
    \item Mayo Clinic. \textit{Hepatitis B: symptoms and causes.} https://www.mayoclinic.org/
\end{itemize}
